% GENERATED FILE. Edit canonical lessons, then run npm run generate:books.
% canonical-source-hash: fnv1a64:8ec22396f64278a9
% canonical-lessons: GE-C327-koecheln, GE-C327-aufkochen, GE-C327-duensten, GE-C327-anbraten, GE-C327-deckel

\chapter{Auf dem Herd}
\label{ch:ge-auf-dem-herd}
\hlchaptermodality{327}

\begin{chapteropening}
\textbf{By the end of this chapter:} \emph{I can describe cooking on the stove: to simmer, to bring to the boil, to steam, to brown quickly and a lid.}

Complete the last lesson of chapter 327: I can describe cooking on the stove: to simmer, to bring to the boil, to steam, to brown quickly and a lid.
\end{chapteropening}

\section[köcheln]{köcheln — to simmer}

\label{lesson:GE-C327-koecheln}



\begin{quote}
Before the new one: say the German for the baking dish, then the German for to brown.
\end{quote}

\subsection*{You'll want to know: köcheln}
\textbf{köcheln} — \textquotedblleft{}to simmer\textquotedblright{}.

\begin{quote}
Die Suppe köchelte eine Stunde.
\end{quote}

\begin{quote}
\emph{The soup simmered for an hour.}
\end{quote}

A small form of \emph{kochen}: gentle cooking.

\begin{cousinweb}
\emph{köcheln} is built on \emph{kochen} \textquotedblleft{}to cook\textquotedblright{}, Old High German \emph{kochōn}, which EtymDB traces to Latin \emph{coquus} \textquotedblleft{}a cook\textquotedblright{}, from \emph{coquere} \textquotedblleft{}to cook\textquotedblright{}; the ending \emph{-eln} makes it small and gentle.
\end{cousinweb}

\begin{grammarlens}[title={What you've built}]
One more word for telling, explaining and describing at B1.
\end{grammarlens}

\subsection*{Your turn}
\begin{itemize}
\raggedright
  \item \emph{Say it:} \emph{köcheln}
  \item \emph{Say it:} \emph{köcheln}, once more
  \item \emph{Say it:} a sentence of your own with \emph{köcheln}
  \item \emph{Recall:} say the German for the baking dish, then the German for to brown, then say \emph{köcheln} again
\end{itemize}

\begin{tcolorbox}[breakable,colback=teal!4,colframe=teal!35!black,title={Before you move on}]
What does köcheln mean? (\textquotedblleft{}To simmer\textquotedblright{}.) Say it once more.
\end{tcolorbox}
\section[aufkochen]{aufkochen — to bring to the boil}

\label{lesson:GE-C327-aufkochen}



\begin{quote}
Before the new one: say the German for to brown, then the German for to simmer.
\end{quote}

\subsection*{You'll want to know: aufkochen}
\textbf{aufkochen} — \textquotedblleft{}to bring to the boil\textquotedblright{}.

\begin{quote}
Sie ließ die Milch kurz aufkochen.
\end{quote}

\begin{quote}
\emph{She brought the milk briefly to the boil.}
\end{quote}

Separable: \emph{etwas aufkochen lassen}.

\begin{cousinweb}
\emph{aufkochen} is built on \emph{kochen} \textquotedblleft{}to cook\textquotedblright{}, Old High German \emph{kochōn}, an early loan from Latin \emph{coquus} \textquotedblleft{}a cook\textquotedblright{} (EtymDB), with \emph{auf} \textquotedblleft{}up\textquotedblright{}.
\end{cousinweb}

\begin{grammarlens}[title={What you've built}]
One more word for telling, explaining and describing at B1.
\end{grammarlens}

\subsection*{Your turn}
\begin{itemize}
\raggedright
  \item \emph{Say it:} \emph{aufkochen}
  \item \emph{Say it:} \emph{aufkochen}, once more
  \item \emph{Say it:} a sentence of your own with \emph{aufkochen}
  \item \emph{Recall:} say the German for to brown, then the German for to simmer, then say \emph{aufkochen} again
\end{itemize}

\begin{tcolorbox}[breakable,colback=teal!4,colframe=teal!35!black,title={Before you move on}]
What does aufkochen mean? (\textquotedblleft{}To bring to the boil\textquotedblright{}.) Say it once more.
\end{tcolorbox}
\section[dünsten]{dünsten — to steam, to braise}

\label{lesson:GE-C327-duensten}



\begin{quote}
Before the new one: say the German for to simmer, then the German for to bring to the boil.
\end{quote}

\subsection*{You'll want to know: dünsten}
\textbf{dünsten} — \textquotedblleft{}to steam, to braise\textquotedblright{}.

\begin{quote}
Er dünstete das Gemüse in wenig Wasser.
\end{quote}

\begin{quote}
\emph{He steamed the vegetables in a little water.}
\end{quote}

Gentle cooking with little liquid.

\begin{cousinweb}
\emph{dünsten} is built on \emph{Dunst} \textquotedblleft{}vapour, steam\textquotedblright{} (EtymDB), an inherited Germanic word.
\end{cousinweb}

\begin{grammarlens}[title={What you've built}]
One more word for telling, explaining and describing at B1.
\end{grammarlens}

\subsection*{Your turn}
\begin{itemize}
\raggedright
  \item \emph{Say it:} \emph{dünsten}
  \item \emph{Say it:} \emph{dünsten}, once more
  \item \emph{Say it:} a sentence of your own with \emph{dünsten}
  \item \emph{Recall:} say the German for to simmer, then the German for to bring to the boil, then say \emph{dünsten} again
\end{itemize}

\begin{tcolorbox}[breakable,colback=teal!4,colframe=teal!35!black,title={Before you move on}]
What does dünsten mean? (\textquotedblleft{}To steam, to braise\textquotedblright{}.) Say it once more.
\end{tcolorbox}
\section[anbraten]{anbraten — to brown quickly, to sear}

\label{lesson:GE-C327-anbraten}



\begin{quote}
Before the new one: say the German for to bring to the boil, then the German for to steam.
\end{quote}

\subsection*{You'll want to know: anbraten}
\textbf{anbraten} — \textquotedblleft{}to brown quickly, to sear\textquotedblright{}.

\begin{quote}
Zuerst briet sie das Fleisch kurz an.
\end{quote}

\begin{quote}
\emph{First she seared the meat briefly.}
\end{quote}

Separable, strong: \emph{brät an}, \emph{briet an}, \emph{hat angebraten}.

\begin{cousinweb}
\emph{anbraten} is built on \emph{braten} \textquotedblleft{}to roast, to fry\textquotedblright{}, Old High German \emph{brātan} (EtymDB), with \emph{an}.
\end{cousinweb}

\begin{grammarlens}[title={What you've built}]
One more word for telling, explaining and describing at B1.
\end{grammarlens}

\subsection*{Your turn}
\begin{itemize}
\raggedright
  \item \emph{Say it:} \emph{anbraten}
  \item \emph{Say it:} \emph{anbraten}, once more
  \item \emph{Say it:} a sentence of your own with \emph{anbraten}
  \item \emph{Recall:} say the German for to bring to the boil, then the German for to steam, then say \emph{anbraten} again
\end{itemize}

\begin{tcolorbox}[breakable,colback=teal!4,colframe=teal!35!black,title={Before you move on}]
What does anbraten mean? (\textquotedblleft{}To brown quickly, to sear\textquotedblright{}.) Say it once more.
\end{tcolorbox}
\section[der Deckel]{der Deckel — the lid}

\label{lesson:GE-C327-deckel}



\begin{quote}
Before the new one: say the German for to steam, then the German for to brown quickly.
\end{quote}

\subsection*{You'll want to know: der Deckel}
\textbf{der Deckel} — \textquotedblleft{}the lid\textquotedblright{}.

\begin{quote}
Er legte den Deckel auf den Topf.
\end{quote}

\begin{quote}
\emph{He put the lid on the pot.}
\end{quote}

Plural: \emph{die Deckel}.

\begin{cousinweb}
\emph{Deckel} is built on \emph{decken} \textquotedblleft{}to cover\textquotedblright{}, from Proto-Germanic \emph{þakjaną}, which EtymDB builds on \emph{þaką} \textquotedblleft{}thatch, roof\textquotedblright{}.
\end{cousinweb}

\begin{grammarlens}[title={What you've built}]
One more word for telling, explaining and describing at B1.
\end{grammarlens}

\subsection*{Your turn}
\begin{itemize}
\raggedright
  \item \emph{Say it:} \emph{der Deckel}
  \item \emph{Say it:} \emph{der Deckel}, once more
  \item \emph{Say it:} a sentence of your own with \emph{der Deckel}
  \item \emph{Recall:} say the German for to steam, then the German for to brown quickly, then say \emph{der Deckel} again
\end{itemize}

\begin{tcolorbox}[breakable,colback=teal!4,colframe=teal!35!black,title={Before you move on}]
What does der Deckel mean? (\textquotedblleft{}The lid\textquotedblright{}.) Say it once more.
\end{tcolorbox}
